\documentclass{article}
\usepackage[utf8]{inputenc}
\usepackage{amsmath}

\begin{document}

\section{Interactive Matrix Resizer 테스트}

이 파일에서 마우스 커서를 행렬 코드 위에 올리면 나타나는 Hover 툴팁을 통해 행렬 구조를 조작해 보세요.

\subsection{1. 기본 행렬 (Basic Matrices)}
아래 행렬들에 마우스를 올리고 [➕ 행 추가] 또는 [➕ 열 추가]를 클릭해 보세요.

\[
\begin{bmatrix}
    1 & 0 \\
    0 & 1
\end{bmatrix}
\]

\[
\begin{pmatrix}
    a & b \\
    c & d
\end{pmatrix}
\]

\subsection{2. 복잡한 셀 내용 (Complex Cell Content)}
분수나 첨자가 들어있어도 구조가 깨지지 않고 열/행이 정확히 추가/삭제되는지 확인해 보세요.
(예: 열 추가 시 \frac{1}{2} 뒤에 &가 올바르게 붙어야 함)

\[
\begin{vmatrix}
    \frac{\alpha}{\beta} & \sum_{i=1}^n x_i \\
    e^{j\omega t} & \int_0^\infty f(x)dx
\end{vmatrix}
\]

\subsection{3. 기타 지원 환경 (Other Environments)}
행렬뿐만 아니라 cases나 align 환경에서도 작동합니다.

\[
f(x) = \begin{cases}
    x & x > 0 \\
    -x & x < 0
\end{cases}
\]

\begin{align}
    2x + 3y &= 7 \\
    x - y &= 1
\end{align}

\subsection{4. 빈 행렬 테스트}
아래의 빈 행렬에서 시작해 보세요.

\[
\begin{bmatrix}
\end{bmatrix}
\]

\end{document}
